% =============================================================================
% 1 INTRODUÇÃO — primeira versão (18/11/2026), escrita com a skill
% estruturacao-introducao-logica: abre com a conclusão principal; problema,
% razão para esperar o efeito, lacuna e justificativa num argumento
% contínuo que leva ao objetivo; objetivos e organização como subseções.
% Guia UECE p. 64: seção numerada, alinhada à esquerda.
% Fontes: docs/briefing.md §1, §6; docs/decisions.md (D-001, D-004, D-008);
% artifacts/reports/results_final_*-tiers.md (números da abertura);
% docs/qualitative_examples.md (Q011: 10.931 x 10.375).
% =============================================================================
\chapter{Introdução}
\label{cap:introducao}

Este trabalho mostra que, em um agente que traduz perguntas em português
para consultas SQL, acrescentar ao esquema do banco regras semânticas
explícitas aumenta a proporção de consultas corretas, enquanto descrições
de coluna, sozinhas, não produzem diferença significativa. No modelo
\arq{qwen3.8-27b}, a acurácia passou de 60,0\% com as descrições de
catálogo para 87,5\% com a camada semântica, diferença significativa após
o ajuste para comparações múltiplas. No \arq{gpt-oss-120b}, passou de
70,0\% para 82,5\%, na mesma direção, sem atingir significância. O ganho
concentra-se em erros que dependem de regras sobre como consultar a base,
e não do significado das colunas.

Agentes baseados em modelos de linguagem de grande escala
(\textit{large language models}, LLMs) passaram a ser usados para
responder perguntas de negócio diretamente sobre bancos de dados
relacionais \pendente{citação}. Nessa tarefa, chamada de tradução de
linguagem natural para SQL (\textit{text-to-SQL}), o modelo recebe a
pergunta e uma descrição do esquema, isto é, das tabelas, colunas e chaves
do banco, e gera a consulta que a responde \pendente{citação}. O problema
é que a pergunta usa o vocabulário do negócio, enquanto o esquema descreve
a estrutura física dos dados, e a passagem de um para o outro raramente é
única.

Uma mesma pergunta admite, em geral, várias consultas que executam sem
erro e devolvem números diferentes. Na base usada neste trabalho, a
pergunta ``Quantos pedidos têm cliente associado?'' pode ser respondida
contando os pedidos com o identificador do cliente preenchido, o que dá
10.931, ou contando os pedidos ligados a um cliente que existe no
cadastro, o que dá 10.375, porque parte das referências aponta para
clientes inexistentes. As duas consultas são válidas, e só uma
corresponde ao que a organização entende por ``cliente associado''. O erro
é silencioso: a consulta executa, e o usuário recebe um número plausível e
errado.

A informação que decide entre essas alternativas não está no esquema
físico. Ela é uma convenção do negócio: qual é o grão de cada tabela, qual
coluna de data define uma métrica, como tratar referências órfãs, quais
colunas devolver. Se o modelo recebe essas convenções explicitamente, é de
se esperar que escolha a consulta correta com mais frequência. Na
indústria de inteligência de negócios, essas convenções são mantidas em
camadas semânticas, ferramentas como dbt Semantic Layer, Databricks Unity
Catalog, LookML e Cube, que definem métricas, dimensões e relacionamentos
sobre o banco físico \pendente{citação}. Um passo intermediário, mais
simples e mais comum, são os catálogos de dados, que guardam apenas uma
descrição textual de cada tabela e coluna.

Essas duas formas de contexto carregam informações de natureza diferente.
Uma descrição de coluna diz o que a coluna significa; uma regra semântica
diz como a base deve ser consultada. Um relato da comunidade Databricks
comparou os três níveis, apenas as tabelas, as tabelas com descrições de
catálogo e as tabelas com métricas e instruções, e encontrou um ganho
pequeno com as descrições e um ganho grande com a camada semântica
\pendente{citação}. Esse resultado, porém, foi obtido em uma plataforma
proprietária, com agente próprio, e não permite saber se o efeito vem do
conteúdo do contexto ou da plataforma. Os \textit{benchmarks} acadêmicos de
\textit{text-to-SQL}, por sua vez, usam em geral esquemas preparados para
avaliação \pendente{citação}, sem as imperfeições de uma base comercial
real, como as referências órfãs do exemplo acima.

Permanece em aberto, portanto, se o efeito da camada semântica se mantém
quando o contexto é apenas texto enviado ao modelo, com modelos de pesos
abertos e uma base comercial real, e qual parte do contexto o produz. A
resposta tem relevância prática: uma organização que já mantém um catálogo
de dados precisa saber se as descrições de coluna bastam para um agente
ou se é preciso escrever regras. Tem também relevância metodológica, porque
a comparação pode ser feita sem plataforma paga. Neste trabalho, o
experimento inteiro usa ferramentas abertas e modelos servidos
gratuitamente, de modo que qualquer pesquisador pode reproduzi-lo, e a
base é um recorte real e descaracterizado de dados comerciais, autorizado
para uso acadêmico (\autoref{sec:fonte-dados}).

A pergunta de pesquisa é, então: em um agente de tradução de linguagem
natural para SQL, enriquecer a descrição técnica do esquema com
descrições de catálogo e, em seguida, com uma camada semântica altera a
proporção de consultas corretas? Para cada par de condições, a hipótese
nula é a de que não há diferença de acurácia; a hipótese alternativa é
bilateral, de modo que o estudo pode detectar tanto melhora quanto piora.

\section{Objetivos}
\label{sec:objetivos}

\subsection{Objetivo geral}

Avaliar o efeito do nível de representação semântica do esquema (técnico,
catálogo e camada semântica) na correção das consultas SQL geradas e
executadas por agentes baseados em modelos de linguagem.

\subsection{Objetivos específicos}

\begin{alineas}
  \item construir um protocolo experimental pareado, contrabalanceado e
        reprodutível sobre uma base comercial real descaracterizada, com
        ferramentas abertas e sem custo financeiro;
  \item comparar a acurácia do agente nas três condições de contexto, em
        dois modelos de pesos abertos de famílias distintas, com testes
        pareados e ajuste para comparações múltiplas;
  \item classificar os erros das consultas geradas e identificar qual
        nível de contexto resolve cada classe de erro;
  \item medir o custo do contexto adicional em latência;
  \item relacionar o contexto semântico usado às camadas semânticas
        industriais e ao resultado relatado pela comunidade Databricks.
\end{alineas}

\section{Organização do trabalho}
\label{sec:organizacao}

A \autoref{cap:fundamentacao} apresenta os conceitos de modelos de
linguagem, tradução de linguagem natural para SQL, representação do
esquema, camada semântica e avaliação estatística pareada. A
\autoref{cap:trabalhos-relacionados} discute os trabalhos relacionados. A
\autoref{cap:metodologia} descreve a base de dados, as condições de
contexto, o agente, os modelos, o banco de perguntas, o critério de
correção e a análise estatística. A \autoref{cap:resultados} apresenta os
resultados, e a \autoref{cap:discussao} os interpreta. A
\autoref{cap:conclusao} resume as conclusões e indica trabalhos futuros.
